% !tp-example Basic usage
\documentclass{article}
\begin{document}
\rawhtml|<tp-cryptarithm equation="SEND + MORE = MONEY" solution="9567 + 1085 = 10652"></tp-cryptarithm>|
\end{document}

% !tp-example Javascript API
\documentclass{article}
\begin{document}
\wcScript[type=module]|
import "/components/cryptarithm/cryptarithm.js";

const example = document.createDocumentFragment();

const element1 = document.createElement("tp-cryptarithm");
element1.setAttribute("equation", "SEND + MORE = MONEY");
element1.setAttribute("solution", "9567 + 1085 = 10652");
example.append(element1);

document.body.append(example);
|
\end{document}

% !tp-example UN + UN + NEUF = ONZE
\documentclass{article}
\begin{document}
\rawhtml|<tp-cryptarithm equation="UN + UN + NEUF = ONZE" solution="81 + 81 + 1987 = 2149"></tp-cryptarithm>|
\end{document}

% !tp-example CINQ + CINQ + VINGT = TRENTE
\documentclass{article}
\begin{document}
\rawhtml|<tp-cryptarithm equation="CINQ + CINQ + VINGT = TRENTE" solution="6483 + 6483 + 94851 = 107817"></tp-cryptarithm>|
\end{document}

% !tp-example ZERO + NEUF + NEUF + DOUZE = TRENTE
\documentclass{article}
\begin{document}
\rawhtml|<tp-cryptarithm equation="ZERO + NEUF + NEUF + DOUZE = TRENTE" solution="9206 + 3257 + 3257 + 86592 = 102312"></tp-cryptarithm>|
\end{document}
